% 参考文献：按正文首次引用顺序编号，GB/T 7714 著录。[3]、[4] 的卷期页码待核对原文。
\begin{thebibliography}{9}
\setlength{\itemsep}{2pt}
\bibitem{gssap} United States Space Force. Geosynchronous Space Situational Awareness Program[EB/OL]. \url{https://www.spoc.spaceforce.mil/About-Us/Fact-Sheets/Display/Article/2381739/geosynchronous-space-situational-awareness-program}.
\bibitem{zhang2024} 张海涛, 王伟林, 张雅声, 等. 地球静止轨道航天器绕飞持续观测任务轨迹规划与控制[J]. 国防科技大学学报, 2024, 46(1): 74-86.
\bibitem{alfriend2006} ALFRIEND K T, LEE D J, CREAMER N G. Optimal servicing of geosynchronous satellites[J]. Journal of Guidance, Control, and Dynamics, 2006, 29(1): 203-206.
\bibitem{zhou2015} ZHOU Y, YAN Y, HUANG X, et al. Mission planning optimization for multiple geosynchronous satellites refueling[J]. Advances in Space Research, 2015, 56(11): 2612-2625.
\bibitem{wang2025} 王峥嵘, 曹晓瑞, 黄喜元, 等. 地球同步轨道多目标飞掠成像服务任务规划方法[J]. 航天控制, 2025, 43(3): 15-23.
\bibitem{kraft1988} KRAFT D. A software package for sequential quadratic programming[R]. Köln: DFVLR, 1988. DFVLR-FB 88-28.
\bibitem{scipy2020} VIRTANEN P, GOMMERS R, OLIPHANT T E, et al. SciPy 1.0: fundamental algorithms for scientific computing in Python[J]. Nature Methods, 2020, 17(3): 261-272.
\end{thebibliography}
